\begin{thema}{Β}
Δίνεται η συνάρτηση $f(x) = x^3 + a x^2 + 3x + 1$.
\begin{erwthma}
\item Η $f$ είναι δύο φορές παραγωγίσιμη ως πολυωνυμική:
\[ f'(x) = 3x^2 + 2ax + 3 \]
\[ f''(x) = 6x + 2a \]
Γνωρίζουμε ότι παρουσιάζει σημείο καμπής στο σημείο με τετμημένη $x_0 = 1$. Επομένως, $f''(1) = 0$.
\[ 6(1) + 2a = 0 \Leftrightarrow 2a = -6 \Leftrightarrow a = -3 \]
Για $a = -3$, ο τύπος είναι $f(x) = x^3 - 3x^2 + 3x + 1$ και η πρώτη της παράγωγος:
\[ f'(x) = 3x^2 - 6x + 3 = 3(x^2 - 2x + 1) = 3(x-1)^2 \]
Ισχύει $f'(x) \ge 0$ για κάθε $x \in \mathbb{R}$, και επιπλέον μηδενίζεται μόνο στο μεμονωμένο σημείο $x=1$. Επομένως, σύμφωνα με το κριτήριο μονοτονίας, η συνάρτηση $f$ είναι πράγματι \textbf{γνησίως αύξουσα} σε όλο το $\mathbb{R}$.

\item *(Σημείωση: Στην εκφώνηση ζητείται να δειχθεί ότι υπάρχει ρίζα $x_0 \in (0, 1)$ για την εξίσωση $f(x_0)=2$, ωστόσο αυτό περιέχει ένα ανακόλουθο).* 
Υπολογίζουμε την τιμή $f(1)$:
\[ f(1) = 1^3 - 3(1)^2 + 3(1) + 1 = 1 - 3 + 3 + 1 = 2 \]
Άρα η εξίσωση $f(x) = 2$ δέχεται ως προφανή λύση το $x=1$. 
Επειδή η συνάρτηση αποδείχθηκε (στο ερώτημα 1) \textbf{γνησίως αύξουσα} στο $\mathbb{R}$, η λύση $x=1$ είναι η \textbf{μοναδική}. Κατά συνέπεια, δεν υπάρχει καμία ρίζα στο ανοικτό διάστημα $(0, 1)$. Το ζητούμενο της άσκησης πιθανόν να αφορούσε κλειστό διάστημα $[0, 1]$ ή διαφορετική τιμή εξίσωσης. Εμείς αποδείξαμε πάντως τη μοναδικότητα.

\item Ζητείται το όριο $\lim_{x\to0} \frac{f(x)-1}{x}$.
\[ \lim_{x\to0} \frac{x^3 - 3x^2 + 3x + 1 - 1}{x} = \lim_{x\to0} \frac{x^3 - 3x^2 + 3x}{x} = \lim_{x\to0} \frac{x(x^2 - 3x + 3)}{x} \]
Αφού $x \ne 0$, απλοποιούμε και βρίσκουμε:
\[ \lim_{x\to0} (x^2 - 3x + 3) = 0 - 0 + 3 = 3 \]
*(Αξίζει να σημειωθεί ότι το όριο αυτό είναι εξ ορισμού η παράγωγος $f'(0) = 3$).*

\item Η δεύτερη παράγωγος για $a=-3$ είναι $f''(x) = 6x - 6 = 6(x-1)$.
- Για $x > 1 \Rightarrow f''(x) > 0$. Η $f$ είναι \textbf{κυρτή} στο $[1, +\infty)$.
- Για $x < 1 \Rightarrow f''(x) < 0$. Η $f$ είναι \textbf{κοίλη} στο $(-\infty, 1]$.
Όπως είχε δοθεί, το $x=1$ είναι τo μοναδικό σημείο καμπής.
Στο σημείο καμπής $A(1, 2)$ η εφαπτομένη εξάγεται ως εξής:
Έχουμε $f'(1) = 3(1-1)^2 = 0$. Η κλίση είναι μηδέν.
Η εξίσωση είναι: $y - f(1) = f'(1)(x-1) \Leftrightarrow y - 2 = 0 \Leftrightarrow y = 2$.
Είναι δηλαδή μία \textbf{οριζόντια εφαπτομένη}.
\end{erwthma}
\end{thema}
